\documentclass{article}
\usepackage[T1]{fontenc}
\usepackage{lmodern}
\usepackage{graphicx}
\usepackage{amsmath,amssymb}
\usepackage[a4paper,margin=2cm]{geometry}
\usepackage[absolute,overlay]{textpos}
\setlength{\TPHorizModule}{1mm}
\setlength{\TPVertModule}{1mm}
\usepackage[indonesian]{babel}
\usepackage{hyperref}
\setlength{\emergencystretch}{2em}
% unit: Halaman 10

\begin{document}

% === Halaman 10 (python/native) ===

10

Prosedur Pembangkitan Kunci

1.

Pilih sembarang bilangan prima p ( p dapat di-share di antara 

anggota kelompok)

2.

Pilih dua buah bilangan acak, g dan x, dengan syarat g < p, g akar 

primitif dari p, dan  2  x  p – 2 

3.    Hitung y = gx mod p  



     (1)



Hasil dari algoritma ini:


- Kunci publik: tripel (y, g, p)


- Kunci privat: pasangan (x, p)

\end{document}
